\chapter{仓库结构就是你的第二大脑}
\label{ch:repo_layout}

研究代码的"混乱"常常不是因为能力不足，而是因为研究天然是并行探索：同一时间你要维护多个假设、多个实现、多个实验入口、多个输出与多个图表。

当这些东西没有结构地堆在一起，你的大脑就会被迫充当索引：
\textbf{哪个脚本还能跑？哪个输出可信？哪个改动影响了指标？}

仓库结构的目的不是好看，而是\textbf{降低认知负担}：让你不用靠回忆就能判断"这段代码是否可删、这份输出是否可信、这条路径是否可复现"。

\section{真实案例：快速堆代码的代价}

在我早期使用 AI 编码助手时，曾经历过这样一次教训：为了快速验证一个想法，我让 Copilot 帮我生成了大量"能跑"的代码——数据加载、模型定义、训练循环、评估脚本等等，几个小时内就搭起了一个看似完整的框架。

\paragraph{初期的"成功"：}代码确实跑通了，实验也出了结果。我兴奋地继续迭代，不断让 AI 帮我添加新功能：数据增强、不同的模型变体、各种评估指标……每次改动都"立竿见影"，代码量迅速膨胀。

\paragraph{崩溃的开始：}两周后，当我需要准备论文的对照实验时，问题开始暴露：
\begin{itemize}
  \item 我分不清哪个脚本是最新的，哪个是废弃的；
  \item 同样的数据加载逻辑被复制粘贴到了 5 个不同文件里，每个版本略有差异；
  \item baseline 和新方法用了不同的评估代码，结果根本不可比；
  \item 想要复现某个"很好的结果"，却找不到当时用的配置和数据版本。
\end{itemize}

\paragraph{推倒重来：}最终，我不得不停下所有新实验，花了整整三天时间\textbf{重写几乎所有代码}。这次重写不是因为 AI 生成的代码有 bug，而是因为\textbf{缺乏结构}：可复用的核心逻辑和一次性实验脚本混在一起，快速试错的代码没有及时清理，输出散落各处难以追溯。

这次经历让我深刻理解：AI 可以帮你快速产出代码，但\textbf{结构必须由人来设计}。如果在最开始就把"稳定"和"探索"分离开，把输出按 run\_id 组织，后续就不会陷入这种混乱。

\paragraph{案例参考：}这类问题并非个例。当你用 AI 快速堆出一个"能跑"的仓库，却没有把可复用代码与一次性实验入口隔离开来，最终常见结局是：每个模块都要重写，AI 生成片段几乎全被替换（见案例~\ref{case:copilot_rewrite_all}、\ref{case:copilot_overdelegation}）。
\textbf{结构是避免这种返工的第一道防线。}

\section{一个可直接照抄的目录结构（研究友好）}
\begin{verbatim}
repo/
  src/                 # 核心库：可复用、可测试、可维护（慢变量）
  experiments/         # 实验入口：一次性 glue 代码（快变量，可丢弃）
  configs/             # 统一配置：yaml/json（可diff，可追溯）
  data/                # 只放指针或小样本；大数据用外部管理
  outputs/             # 运行产物：按 run_id 组织（可清理/可归档）
  reports/             # 论文图表与结论：从 outputs 自动生成
  scripts/             # 数据准备/下载/评估等工具脚本
  tests/               # 单测 + smoke test（守住底线）
  Makefile             # 常用入口：train/eval/reproduce/test
  README.md
  CLAUDE.md            # AI 编码规则（建议）
\end{verbatim}

\section{快变量与慢变量：把"稳定"从"探索"里分离出来}
建议把仓库里的东西分成两类：
\begin{itemize}
  \item \textbf{慢变量（stable）：}会被长期维护、反复复用、需要测试覆盖的部分。
  \item \textbf{快变量（explore）：}为某个假设快速试错的入口脚本、短命 glue、一次性分析。
\end{itemize}

在本书的术语里：\texttt{src/} 是慢变量，\texttt{experiments/} 是快变量。

\textbf{经验法则：}探索可以脏，但核心库必须干净；探索可以快，但评估必须稳定。

\paragraph{为什么这个分离如此重要？}在我重写代码的经历中，最大的困扰就是\textbf{无法区分资产和消耗品}。当所有代码都混在一起时，你不敢删除任何东西（怕删错重要功能），也不敢大改（怕影响到其他实验）。一旦明确了 \texttt{src/} 是资产、\texttt{experiments/} 是消耗品，心理负担就会大大减轻：
\begin{itemize}
  \item 对 \texttt{src/} 的改动要谨慎，需要测试；
  \item 对 \texttt{experiments/} 的改动可以随意，试完就删。
\end{itemize}

\section{每个目录的 Definition of Done（DoD）}
目录名只有在"什么该放、什么不该放"足够清晰时，才会真正降低混乱。

\subsection{src/：核心库（可复用、可测试）}
\begin{itemize}
  \item 放可复用模块：数据加载、模型组件、损失、评估、通用工具。
  \item 必须可测试：至少有 smoke test 覆盖关键链路。
  \item 禁止硬编码：不要出现只对某个 run 有用的路径/参数。
\end{itemize}

\paragraph{反例：}在我的重写案例中，原本的"数据加载"代码硬编码了某个特定实验的路径和预处理方式，导致新实验不得不复制粘贴并修改。如果当初就把路径和参数作为函数参数传入，就不会有这个问题。

\subsection{experiments/：实验入口（可丢弃）}
\begin{itemize}
  \item 只放入口与 glue：\textbf{允许短命}，用完即删。
  \item 任何被证明有价值、会复用的逻辑，一旦稳定就迁移到 \texttt{src/}。
\end{itemize}

\paragraph{实践建议：}每个实验脚本以日期或 run\_id 命名，比如 \texttt{2026-02-01\_baseline.py}。这样你一眼就能看出哪些是老实验，哪些是新实验。定期（比如每周）清理超过一个月且没有价值的脚本。

\subsection{configs/：配置（可追溯）}
\begin{itemize}
  \item 每个"论文候选结论"都必须对应一个 config（或其可追溯的生成方式）。
  \item config 必须能展开成最终参数（避免默认值漂移）。
\end{itemize}

\subsection{outputs/：产物（可清理/可归档）}
\begin{itemize}
  \item 只放运行产物，并且\textbf{按 run\_id} 组织。
  \item 禁止覆盖：同一个 run\_id 目录不复用/不手工改。
  \item 重要产物做归档：进入长期存储，仓库不背大二进制。
\end{itemize}

\subsection{reports/：图表与结论（可再生成）}
\begin{itemize}
  \item 尽量只放脚本生成的图表/表格与结论草稿。
  \item 论文图表必须能从 \texttt{outputs/} 重新生成，避免手工拖拽。
\end{itemize}

\subsection{tests/：测试（守住底线）}
\begin{itemize}
  \item 至少一个 1--3 分钟的 smoke test：跑通数据加载 $\rightarrow$ 前向 $\rightarrow$ 损失 $\rightarrow$ 评估。
  \item 关键函数做断言：shape、NaN、范围、数据泄漏信号等。
\end{itemize}

\section{入口收敛：让"怎么跑"变得显眼}
研究里最常见的浪费是：别人（包括未来的你）不知道怎么跑。

建议把所有常用入口收敛到 \texttt{Makefile}（或等价任务工具）：
\begin{itemize}
  \item \texttt{make test}
  \item \texttt{make train CONFIG=...}
  \item \texttt{make eval RUN=...}
  \item \texttt{make reproduce RUN=...}
\end{itemize}

当入口足够少、足够稳定时，你就能把复杂性关在内部，把可复现性暴露在外部。

\section{渐进式重构：从混乱项目迁移到良好结构}

如果你已经有一个"混乱"的老项目，不要试图一次性推倒重来。以下是一个循序渐进的重构流程：

\subsection{第一步：识别并区分慢变量与快变量}

\begin{enumerate}
  \item \textbf{审视现有代码：}扫描整个项目，标记出哪些模块是核心功能（会长期维护、反复使用），哪些是一次性试验脚本。

  \item \textbf{新建目录：}
  \begin{itemize}
    \item 创建 \texttt{src/}，把确定会复用的模块迁移进去；
    \item 创建 \texttt{experiments/}，将各种杂散的运行脚本和临时代码移入其中。
  \end{itemize}

  \item \textbf{完成粗分层：}这一步完成后，项目结构会开始变清晰，你脑中的索引负担也随之减轻。不要求完美，先把明显的部分分开即可。
\end{enumerate}

\subsection{第二步：提取配置与参数}

\begin{enumerate}
  \item \textbf{找出硬编码：}扫描代码，找出所有硬编码的关键参数（学习率、batch size、路径等）。

  \item \textbf{创建配置文件：}在 \texttt{configs/} 目录下创建 YAML 或 JSON 文件，把参数集中管理。

  \item \textbf{逐步替换：}不要一次性改完所有代码，而是每次重构一个脚本，确保它能正常运行后再继续下一个。
\end{enumerate}

\paragraph{实践经验：}在我的重构过程中，光是提取配置这一步就帮我发现了 3 个隐藏的 bug——原来不同脚本里"看似相同"的参数其实值不一样，导致结果不可比。

\subsection{第三步：统一输出路径}

\begin{enumerate}
  \item \textbf{设定规范：}决定使用 \texttt{outputs/<run\_id>/} 的结构。

  \item \textbf{修改代码：}调整所有实验入口，让输出都按照 run\_id 归档。

  \item \textbf{清理旧输出：}将散落的旧输出文件整理归档或删除，保持 outputs 目录干净。
\end{enumerate}

\subsection{第四步：添加基本测试}

\begin{enumerate}
  \item \textbf{编写 smoke test：}创建一个 1-3 分钟的快速测试，验证核心流程能跑通。

  \item \textbf{每次重构后运行：}确保改动没有破坏基本功能。

  \item \textbf{逐步增加覆盖：}随着重构进展，逐步为关键函数添加单元测试。
\end{enumerate}

\paragraph{关键原则：}循序渐进很重要。不要试图一次性做完全部重构，而是每步都确保现有功能正常再进行下一步。我的重写用了三天，但如果当初采用渐进式重构，可能一周内分散进行，不会影响正常实验进度。

\section{多任务/多项目的目录层级管理}

当面对多个相关但独立的研究任务时，如何组织目录是一个重要问题。

\subsection{原则：优先考虑独立仓库}

\textbf{最佳实践：}每个研究项目使用独立的代码仓库。这样可以：
\begin{itemize}
  \item 避免不同项目的依赖冲突；
  \item 确保版本控制的独立性；
  \item 简化复现流程（每个项目有独立的环境和依赖）。
\end{itemize}

\paragraph{经验法则：}如果两个项目的依赖版本、数据集、运行环境差异较大，\textbf{强烈建议分仓库}。

\subsection{单仓库多项目的分层结构}

如果确实需要在一个仓库中管理多个相关任务（比如同一论文的多个实验、或者共享大量代码的子项目），可以采用以下结构：

\begin{verbatim}
repo/
  projectA/
    src/
    experiments/
    configs/
    outputs/
    README.md
  projectB/
    src/
    experiments/
    configs/
    outputs/
    README.md
  common/
    src/               # 跨项目共享的核心代码
    tests/             # 共享功能的测试
    scripts/           # 通用工具脚本
  README.md            # 总体说明
  Makefile             # 跨项目的通用命令
\end{verbatim}

\subsection{命名约定和环境隔离}

\begin{itemize}
  \item \textbf{配置文件命名：}使用项目前缀，如 \texttt{configs/projectA\_baseline.yaml}、\texttt{configs/projectB\_ablation.yaml}。

  \item \textbf{输出目录：}可以在项目子目录下各自维护 outputs，或者统一在根目录但使用项目名作为前缀：\texttt{outputs/projectA/<run\_id>/}。

  \item \textbf{环境管理：}即使代码在一个仓库，也建议为不同项目维护独立的虚拟环境或 Docker 容器，防止依赖冲突。
\end{itemize}

\subsection{何时共享代码，何时复制代码}

\textbf{共享代码到 common/ 的时机：}
\begin{itemize}
  \item 多个项目都需要使用的基础工具函数；
  \item 通用的数据加载或预处理逻辑；
  \item 标准化的评估指标计算。
\end{itemize}

\textbf{复制代码的时机：}
\begin{itemize}
  \item 代码还在快速变化，不同项目需求可能分化；
  \item 共享会导致过度耦合，影响独立性；
  \item 项目即将结束，未来不再需要同步更新。
\end{itemize}

\paragraph{实践建议：}开始时宁可复制，等到代码真正稳定、确认需要共享时再提取到 common/。过早抽象会导致频繁修改共享代码，反而增加维护负担。

\section{10 分钟动作：把当前项目"分层一次"}

如果你现在只做一件事：把当前仓库粗略分成慢变量与快变量。

\begin{enumerate}
  \item 新建 \texttt{src/}，把确定会复用的模块迁进去。
  \item 新建 \texttt{experiments/}，把所有入口脚本搬进去，允许它们短命。
  \item 新建 \texttt{configs/}，把关键参数从脚本里抽出来。
  \item 规定输出统一进 \texttt{outputs/<run\_id>/}。
\end{enumerate}

你会立刻感觉"可控性"提升：因为你开始能区分什么是资产、什么是一次性消耗品。

\paragraph{从个人经验看：}如果我当初在开始使用 AI 编码助手时就建立了这个结构，那三天的重写工作完全可以避免。良好的结构不是为了好看，而是为了\textbf{让你的大脑不再被迫记住一切细节}，让仓库本身成为你可靠的"第二大脑"。
